% !TEX root = task-6-2-3.tex
% Editable task bank. Append new BankTask blocks at the end of the document.
\documentclass[12pt,a4paper]{article}
\usepackage{physics-textbook}
\usepackage{graphicx,needspace}
\usetikzlibrary{decorations.pathmorphing}
\newif\ifShowSolutions
\ifdefined\TutorEdition
  \ShowSolutionstrue
\else
  \ShowSolutionsfalse % Change to \ShowSolutionstrue to display solutions.
\fi
\newcounter{banktask}
\newsavebox{\BankTaskBox}
\newcommand{\BankTaskContent}[3]{%
  \noindent\textbf{\thebanktask.}\quad #2\par
  \ifShowSolutions
    \Needspace{50mm}% Keep the solution heading with its opening text or diagram.
    #3\par
    \ifdefined\TutorEdition
      \ifcsname TutorSource#1\endcsname
        \nopagebreak\smallskip{\small\csname TutorSource#1\endcsname\par}%
      \fi
    \fi
  \fi
}
\newcommand{\BankTask}[3]{%
  \refstepcounter{banktask}\par\medskip
  \sbox{\BankTaskBox}{\begin{minipage}{\linewidth}%
    \BankTaskContent{#1}{#2}{#3}\end{minipage}}%
  \ifdim\dimexpr\ht\BankTaskBox+\dp\BankTaskBox\relax>\dimexpr\textheight-20mm\relax
    \Needspace{35mm}\BankTaskContent{#1}{#2}{#3}%
  \else
    \Needspace{\dimexpr\ht\BankTaskBox+\dp\BankTaskBox\relax}%
    \noindent\usebox{\BankTaskBox}\par
  \fi
  \medskip
}
\begin{document}
\BookSubsection{6.2.3}{Давление света}


% Source sdamgia-19738: https://phys-ege.sdamgia.ru/problem?id=19738
\BankTask{sdamgia-19738}{
\TaskMatchingColumns{Физический опыт}{А) При освещении ярким светом металлической пластины конденсатора из неё вылетают электроны — это можно зарегистрировать, включив конденсатор в электрическую цепь.\\Б) Если поместить внутрь тщательно вакуумированной колбы лёгкую крыльчатку и направить на неё яркий свет, то крыльчатка будет вращаться.}{Физическое явление}{1) давление света\\2) преломление света\\3) фотоэффект\\4) интерференция света}
\begin{center}\begin{tabular}{|c|c|}\hline А&Б\\\hline &\\\hline\end{tabular}\end{center}
}{
\smallskip\noindent\textbf{Решение}\par
А.  При освещении светом металлической пластины наблюдается явление фотоэффекта.
Б.  Крыльчатка вращается в результате давления, которое оказывает свет.

Ответ: $31$.
}

% Source sdamgia-3084: https://phys-ege.sdamgia.ru/problem?id=3084
\BankTask{sdamgia-3084}{
Источник в монохроматическом пучке параллельных лучей за время $\Delta t=8 \cdot 10 ^{ - 4 } c$ излучает $N=5 \cdot 10 ^{ 14 }$ фотонов. Лучи падают по нормали на площадку $S=0{,}7\,\text{см} ^2$ и создают давление $P=1{,}5 \cdot 10 ^{ - 5 } \,\text{Па}$ . При этом $40\%$ фотонов отражается, а $60\%$ поглощается. Определите длину волны излучения.
}{
\smallskip\noindent\textbf{Решение}\par
Выражение для давления света 

 $P=P_{\,\text{отр}} + P_{\,\text{погл}}= \frac{N_{\,\text{отр}}\Delta p_{\,\text{отр}} + N_{\,\text{погл}}\Delta p_{\,\text{погл}}}{S\Delta t} .$   ($1$) 

 Формула ($1$) следует из: $\vec F= \frac{\Delta \vec p}{\Delta t}$ и $P= \frac{F}{S} .$
Формулы для изменения импульса фотона при отражении и поглощении лучей $\Delta p_{\,\text{отр}}=2p,$ $\Delta p_{\,\text{погл}}=p,$ число отраженных $N_{\,\text{отр}}=0{,}4N$ и поглощенных $N_{\,\text{погл}}=0{,}6N$ фотонов. 

 Тогда выражение ($1$) принимает вид $P= \frac{1{,}4Np}{S\Delta t} .$ 

 Для импульса фотона $p= \frac{h}{\lambda} .$ 

 Выражение для длины волны излучения $\lambda = \frac{1{,}4Nh}{PS\Delta t}$

Ответ: $\lambda = \frac{1{,}4 \cdot 5 \cdot 10 ^{ 14 } \cdot 6{,}6 \cdot 10 ^{ - 34 }}{1{,}5 \cdot 10 ^{ - 5 } \cdot 0{,}7 \cdot 10 ^{ - 4 } \cdot 8 \cdot 10 ^{ - 4 }} =5{,}5 \cdot 10 ^{ - 7 } \,\text{м}.$
}

% Source sdamgia-6365: https://phys-ege.sdamgia.ru/problem?id=6365
\BankTask{sdamgia-6365}{
Давление света от Солнца, который падает перпендикулярно на абсолютно черную поверхность, на орбите Земли составляет около p  =  $5$$\cdot$$10$$^{-6}$ Па. Оцените концентрацию n фотонов в солнечном излучении, считая, что все они имеют длину волны $\lambda$  =  $500\,\text{нм}$.
}{
\smallskip\noindent\textbf{Решение}\par
Каждый фотон имеет скорость c и импульс $P= \frac{h}{\lambda} .$ За время $\Delta t$ на площадку S попадают фотоны из объема $Sc\Delta t.$ В этом объеме содержится $nSc\Delta t$ фотонов с общим импульсом $PnSc\Delta t.$ Фотоны при поглощении абсолютно черной поверхностью полностью передают ей свой импульс. Следовательно, сила давления на площадку S равна $F = \frac{PnSc\Delta t}{\Delta t} = PncS = \frac{hncS}{\lambda} .$ Таким образом, давление равно $p = \frac{F}{S} = \frac{hnc}{\lambda}$ Откуда: 

\[n = \frac{p \lambda}{hc} = \frac{5 \cdot 10 ^{ - 6 } \cdot 5 \cdot 10 ^{ - 7 }}{6{,}6 \cdot 10 ^{ - 34 } \cdot 3 \cdot 10 ^{8}} \approx1{,}3 \cdot 10 ^{ 13 } \,\text{м} ^{ - 3 } .\]

Ответ: $1{,}3$·$10$$^{13}$ м$^{-3}$.
}

% Source sdamgia-6471: https://phys-ege.sdamgia.ru/problem?id=6471
\BankTask{sdamgia-6471}{
Лазер испускает световой импульс с энергией W  =  $12\,\text{Дж}$. Свет от лазера падает перпендикулярно на плоское зеркало площадью S  =  $10\,\text{см}$$^{2}$. Определите длительность импульса $\tau$, если среднее давление света на зеркало равно p  =  $1$ кПа.
}{
\smallskip\noindent\textbf{Решение}\par
$1$.  При отражении от зеркала импульс фотонов изменяется на $\Delta p_{\,\text{ф}}= \frac{2h\nu}{c} ,$ где $\nu$ — частота света, излучаемого лазером; h  — постоянная Планка; с  — скорость света в вакууме. 

$2$.  Пусть N  — число фотонов, излученных лазером, тогда W  =  Nh$\nu$. 

$3$.  Модуль силы, действующей на зеркало $F= \frac{N\Delta p_{\,\text{ф}}}{\tau} .$ Давление света на зеркало: 

\[P= \frac{F}{S} = \frac{N\Delta p_{\,\text{ф}}}{S\tau} = \frac{2Nh v}{cS\tau} = \frac{2W}{cS\tau} ,\]

откуда находим длительность импульса света:

\[\tau= \frac{2W}{cSP} = \frac{24}{3 \cdot 10 ^{8} \cdot 10 ^{ - 3 } \cdot 10 ^3} =8 \cdot 10 ^{ - 8 } =80\,\text{нс}.\]

Ответ: $\tau=80\,\text{нс}.$
}

% Source sdamgia-8962: https://phys-ege.sdamgia.ru/problem?id=8962
\BankTask{sdamgia-8962}{
Параллельный пучок света с длиной волны $\lambda$ = $600\,\text{нм}$ и концентрацией фотонов n = $10$$^{14}$ м$^{-3}$ нормально падает на идеальное зеркало, равномерно освещая всю его поверхность, площадь которой равна S = $1\,\text{м}$$^{2}$. Чему равен модуль силы F давления этого светового пучка на зеркало?
}{
\smallskip\noindent\textbf{Решение}\par
$1$.  Каждый фотон в пучке света имеет энергию $h\nu$ и импульс $h\nu/c.$ Поскольку $\nu = c/\lambda,$ импульс фотона равен $h/\lambda.$
$2$.  При нормальном падении света на идеальное зеркало знак импульса меняется на противоположный, так что каждый фотон передает зеркалу импульс $2h/\lambda.$
$3$.  По второму закону Ньютона сила равна скорости изменения импульса p тела: $F = \Delta p/\Delta t.$
$4$.  За единицу времени ($\Delta t$ = $1\,\text{с}$) от зеркала отразятся все фотоны, находящиеся в цилиндре с высотой, равной скорости света с и площадью основания S, и движущиеся в направлении зеркала, так что $\Delta p = n \cdot cS \cdot 2h/\lambda\Delta t,$ и 

\[F= \frac{2nhcS}{\lambda} =6{,}6 \cdot 10 ^{ - 5 } \,\text{Н}.\]

Ответ: $F= \frac{2nhcS}{\lambda} =6{,}6 \cdot 10 ^{ - 5 } \,\text{Н}.$
}

% Source sdamgia-11687: https://phys-ege.sdamgia.ru/problem?id=11687
\BankTask{sdamgia-11687}{
Для межпланетных полетов в космосе предлагают использовать «солнечный парус»  — большое зеркало, расположенное перпендикулярно солнечным лучам. При их отражении от этого зеркала возникает сила в направлении падающих лучей, которая может ускорять космический корабль. Оцените эту силу F при следующих предположениях: площадь полностью отражающего свет зеркала равна S = $1000\,\text{м}$$^{2}$, а солнечная постоянная в месте нахождения корабля с зеркалом C = $1{,}5$ кВт/м$^{2}$. Солнечная постоянная  — это энергия фотонов, падающих в единицу времени на единицу площади поверхности, перпендикулярной лучам света от Солнца.
}{
\smallskip\noindent\textbf{Решение}\par
$1$.  Согласно определению солнечной постоянной, $C = nEc,$ где n  — концентрация фотонов, c  — скорость света в вакууме, а E  — энергия одного фотона.
$2$.  Поскольку связь энергии E и импульса p фотона имеет вид $E = pc,$ то $C = npc \cdot c,$ откуда плотность потока импульса фотонов (то есть импульс, приносимый фотонами за единицу времени к единице площади) равна $npc = \frac{C}{c} .$
$3$.  Согласно второму закону Ньютона, сила равна скорости изменения импульса всех фотонов, падающих за единицу времени на площадь S, то есть npc · S.
$4$.  Поскольку зеркало полностью отражает свет, то импульс каждого фотона изменяется на $2$p, и $F = 2CS/c = 0{,}01\,\text{Н}.$

Ответ: $F = 2CS/c = 1 \cdot 10 ^{ - 2 } \,\text{Н}.$
}

% Source sdamgia-33055: https://phys-ege.sdamgia.ru/problem?id=33055
\BankTask{sdamgia-33055}{
На идеальное зеркало перпендикулярно ему падают лучи монохроматического света. При их отражении возникает сила давления F, действующая на зеркало. Найти F, если мощность падающего света Р  =  $300$ кВт.
}{
\smallskip\noindent\textbf{Решение}\par
При зеркальном отражении света происходит изменение направления светового пучка, при этом изменение импульса равно $\Delta p.$ Это изменение можно найти из второго закона Ньютона в импульсной форме

\[\Delta p=F\Delta t.\]

Мощность падающего света равна энергии падающих фотонов в единицу времени 

\[P= \frac{Nh\nu}{\Delta t} .\]

Поскольку импульс светового пучка изменился на противоположный, причем $\vec p_1= - \vec p_2,$ то изменение импульса равно $\Delta p=2p_1.$
Импульс фотона $p_1= \frac{h\nu}{c} .$ Тогда сила равна 

\[F= \frac{2h\nu N}{c\Delta t} .\]

С учетом формулы мощности получаем:

\[F= \frac{2P}{c} = \frac{2 \cdot 3 \cdot 10 ^{5}}{3 \cdot 10 ^{8}} =2\,\text{мН}.\]

Ответ: $2$ мН.
}

% Source sdamgia-6436: https://phys-ege.sdamgia.ru/problem?id=6436
\BankTask{sdamgia-6436}{
Лазер испускает световой импульс с энергией W  =  $3\,\text{Дж}$ и длительностью $\tau$ = $10\,\text{нс}$. Свет от лазера падает перпендикулярно на плоское зеркало площадью S  =  $10\,\text{см}$$^{2}$. Какое среднее давление окажет свет на зеркало?
}{
\smallskip\noindent\textbf{Решение}\par
$1$.  При отражении от зеркала импульс фотонов изменяется на $\Delta p_{\,\text{ф}}= \frac{2h\nu}{c} ,$ где $\nu$ — частота света, излучаемого лазером; h  — постоянная Планка; с  — скорость света в вакууме. 
$2$.  Пусть N  — число фотонов, излученных лазером, тогда W  =  Nh$\nu$. 
$3$.  Модуль силы, действующей на зеркало $F= \frac{N\Delta p_{\,\text{ф}}}{\tau} .$ Давление света на зеркало: 

\[P= \frac{F}{S} = \frac{N\Delta p_{\,\text{ф}}}{S\tau} = \frac{2Nh v}{cS\tau} = \frac{2W}{cS\tau} = \frac{2 \cdot 3}{3 \cdot 10 ^{8} \cdot 10 ^{ - 3 } \cdot 10 ^{ - 8 }} =2 \cdot 10 ^3 \,\text{Па}.\]

Ответ: $P=2 \cdot 10 ^3 \,\text{Па}.$
}

% Source sdamgia-25425: https://phys-ege.sdamgia.ru/problem?id=25425
\BankTask{sdamgia-25425}{
Монохроматический пучок света с длиной волны $1{,}1 \cdot 10 ^{ - 10 }$ метров падает на пластину и оказывает давление $1{,}26 \cdot 10 ^{ - 6 } \,\text{Па}.$ $70$\% света отражается, остальное пластинка пропускает. Найти концентрацию фотонов в пучке света. Считать, что фотоны в пучке расположены равномерно.
}{
\smallskip\noindent\textbf{Решение}\par
На пластину могут оказывать давление те фотоны, которые отражаются, то есть $0{,}7N.$ При отражении фотоны меняют направление своего движения на противоположное, поэтому изменение импульса одного фотона равно $\Delta p_{\,\text{ф}}=2p_{\,\text{ф}}.$ Импульс фотона $p_{\,\text{ф}}= \frac{h}{\lambda} .$ Второй закон Ньютона в импульсной форме $F= \frac{0{,}7N\Delta p_{\,\text{ф}}}{\Delta t} .$ Поэтому давление, которое оказывают отраженные от пластины фотоны равно $p= \frac{F}{S} .$ За промежуток времени $\Delta t$ объем пучка света равен $V=Sc\Delta t.$ Объединяя все формулы, получаем $p= \frac{1{,}4 Nhc}{\lambda V} .$
Так как по определению концентрация фотонов $n= \frac{N}{V} ,$ находим ее значение:

\[n= \frac{p\lambda}{1{,}4hc} = \frac{1{,}26 \cdot 10 ^{ - 6 } \cdot 1{,}1 \cdot 10 ^{ - 10 }}{1{,}4 \cdot 6{,}6 \cdot 10 ^{ - 34 } \cdot 3 \cdot 10 ^{8}} =5 \cdot 10 ^{8} \frac{1}{\,\text{м} ^3} .\]

Ответ: $5 \cdot 10 ^{8} \frac{1}{\,\text{м} ^3} .$
}

% Source sdamgia-25429: https://phys-ege.sdamgia.ru/problem?id=25429
\BankTask{sdamgia-25429}{
На пластину падает тонкий пучок света и оказывает на нее давление $P=2 \cdot 10 ^{ - 5 } \,\text{Па},$ концентрация фотонов $n=10 ^{14}\,\text{м}^{-3},$ $80$\% фотонов поглощаются, остальные пролетают насквозь. Найти частоту света.
}{
\smallskip\noindent\textbf{Решение}\par
Давление на пластину оказывают те фотоны, которые поглощаются. Таких фотонов по условию $0{,}8N.$ Переданный пластине импульс равен суммарному импульсу поглощенных фотонов. Импульс одного фотона $p_1= \frac{h v}{c} .$
Из второго закон Ньютона в импульсной форме сила давления фотонов на пластину равна $F= \frac{0{,}8Np_1}{\Delta t} .$ Давление, которое оказывает пучок фотонов на пластину $p= \frac{F}{S} .$ За время $\Delta t$ объем пучка фотонов будет $V=Sc\Delta t.$ Тогда давление фотонов $p= \frac{0{,}8Nh v}{V} .$ Учитывая, что концентрация $n= \frac{N}{V} ,$ находим частоту света:

\[v = \frac{p}{0{,}8nh} = \frac{2 \cdot 10 ^{ - 5 }}{0{,}8 \cdot 10 ^{ 14 } \cdot 6{,}6 \cdot 10 ^{ - 34 }} \approx 3{,}8 \cdot 10 ^{ 14 } \,\text{Гц}.\]

Ответ: $3{,}8 \cdot 10 ^{ 14 } \,\text{Гц}.$
}

\end{document}
